% Thorne & Campolattaro 1967, ApJ 149, 591
% Page 20 (p. 610) transcription

radial perturbations (see Bardeen, Thorne, and Meltzer 1966), so also here, the precise analytic
conditions which these demands place on the eigenfunctions depend upon the distribution of
pressure, density, and adiabatic index near the surface of the equilibrium configuration. In all
cases the eigenfunctions will have to assume the regular form (15b) at the surface, but the
relation of the constants $k_j$, $h_j$, $W_j$, and $v_j$ to each other will vary from case to case.

All cases of current interest fall into two classes: (1) \textit{Absolute zero temperature at the
surface}; example---the Harrison-Wakano-Wheeler configurations (see Harrison et al.\ 1965). In
this case the pressure, density, and adiabatic index near the surface have the form
%
\begin{equation}
p \sim p_0(R - r) \,, \qquad
\rho \sim p_0 + p_1(R - r) \,, \qquad
\gamma p \sim \gamma_0 + \gamma_1(R - r) \,;
\tag{D4}
\end{equation}
%
and, consequently, all five solutions to the eigenequations (14) have the regular form (15b) near
the surface. However, only four of the five solutions are acceptable because only four have
vanishing Lagrangian change in pressure.

(2) \textit{Polytropic pressure-density relation near the surface}; example---any hot stellar model.
In this case the pressure, density, and adiabatic index near the surface are related by
%
\begin{equation}
p = Q\rho^{1+1/N} \,, \qquad \gamma = \gamma_0 \,;
\tag{D5}
\end{equation}
%
and the equation of hydrostatic equilibrium then sets up the pressure and density distributions
%
\begin{equation}
p = a(R - r)^{N+1} \,, \qquad \rho = b(R - r)^N \,.
\tag{D6}
\end{equation}
%
Here $a$, $b$, $\gamma_0$, and $Q$ are constants. An examination of the eigenequations (14) near the
surface for such configurations reveals that only four of the five solutions have the regular form
(15b). These four solutions are physically acceptable because equations (D5) and (D6) guarantee
that $\gamma p$ vanishes at the surface, and therefore---for regular solutions---that $\Delta p$ vanishes
(see eq.\ [16]). The fifth, unacceptable, solution has the surface behavior
%
\begin{align}
K &= k_0 + k_1(R - r) + \cdots \,, \quad
H = h_0 + h_1(R - r) + \cdots \,, \quad
W = G(R - r)^{-N} \,, \notag \\[4pt]
V &= -\frac{G(1 - 2M/R)^{1/2}\,(N\gamma_0 - N - 1)}{\omega^2 R}\,(R - r)^{-N} \,.
\tag{D7}
\end{align}
%
Here $k_0$, $k_1$, $h_0$, $h_1$ are constants freely adjustable by the addition of varying amounts of
the regular solutions (15b); while $G$ is the arbitrary constant which characterizes the solution.

\section*{APPENDIX E}
\subsection*{BEHAVIOR OF THE EVEN-PARITY EIGENFUNCTIONS AT $r = \infty$}

Outside the star's surface ($r > R$) the eigenequations (14b) and (14c) become vacuous, and
(14a) and (14b) alone govern the behavior of the gravitational-wave functions $H_0$ and $K$.
Equations (14a) and (14b) form a third-order differential system and have three independent
solutions. However, only one of these solutions is physically acceptable since only one joins
smoothly ($K'$, $K$, and $H$ continuous) at the star's surface to the perturbed interior geometry.

In the wave zone ($r \gg R$, $r \gg 1/|\omega|$) the eigenequations (14a) and (14b) take on the
simplified form
%
\begin{align}
H' + \bigl[1 + l(l+1)/2\bigr] H/r
  &= K'' + 3K'/r + \bigl[1 - l(l+1)/2\bigr] K/r^2 \,,
\tag{E1a}\\[4pt]
2H/r^2
  &= K'' + 2K'/r + 2\bigl[1 - l(l+1)/2\bigr] K/r^2 + \omega^2\bigl[1 + 4M/r\bigr] K \,.
\tag{E1b}
\end{align}
%
By combining equations (E1a), (E1b), and $d(\text{E1b})/dr$, one can eliminate $H$:
%
\begin{equation}
\bigl[K'' + \omega^2(1 + 4M/r)\,K\bigr]'
+ \bigl[3 + l(l+1)/2\bigr]\,r^{-1}
\times \bigl[K'' + \omega^2(1 + 4M/r)\,K\bigr] = 0 \,.
\tag{E2}
\end{equation}
